\section{Lunar Lander}

\subsection*{(a)}
Code uses a Monte Carlo return-to-go estimate of the advantage $\delta_t^\pi = G_t - V_w(x_t)$,
where \(G_t\) is the discounted return starting at time \(t\). This sits on the
low-bias, high-variance side of the tradeoff. It avoids bootstrapping from the
critic, but it depends on sampled episode returns, which can be noisy.

\subsection*{(b)}
Standardizing the returns keeps the targets on a more reasonable numerical
scale. Raw returns can vary a lot across episodes, so without standardization
the actor and critic updates can have gradients that are too large or too small.
This usually makes training more stable while keeping the same basic learning
signal.

\subsection*{(c)}
Using the function in the actor loss makes the advantage act
like a fixed weight in the policy-gradient update. This is important because the
actor update should change the policy, not backpropagate through the critic or
the return estimate. Without stopping the gradient, the actor loss could also
change the value estimate in a way that is not the intended critic update.

\subsection*{(d)}
Another option would be a model-based controller. With a known or learned
dynamics model for the lander, one could solve a trajectory optimization or MPC
problem with constraints for fuel, thrust limits, and safe landing conditions.
This would use more assumptions than model-free RL, but it could be more sample
efficient and easier to make safe. The main drawback is that performance depends
heavily on the model and cost design being accurate.
